% Copy-ready LaTeX fragment for Questions 1, 3(a), and 4(a).
% Requires \usepackage{amsmath,amssymb} in the main document preamble.

\section*{Question 1}

\subsection*{(a)}

Yes, markets are complete. Let the payoff matrix be
\[
X=
\begin{pmatrix}
10 & 10 & 0 \\
20 & 10 & 10 \\
30 & 10 & 10
\end{pmatrix},
\]
where each column contains the payoff of one asset across the three states.
Markets are complete if and only if there is a unique matrix $\Omega$ such that
$X\Omega=I$. In this case,
\[
\Omega=X^{-1}
=
\begin{pmatrix}
0 & -\frac{1}{10} & \frac{1}{10} \\
\frac{1}{10} & \frac{1}{10} & -\frac{1}{10} \\
-\frac{1}{10} & \frac{1}{5} & -\frac{1}{10}
\end{pmatrix},
\qquad
X\Omega=I.
\]
Because this inverse exists and is unique, each contingent claim can be
replicated by a unique portfolio of the three assets. Therefore, markets are
complete.

\subsection*{(b)}

The identity matrix is the payoff matrix of the three contingent claims. Since
$X\Omega=I$, the law of one price implies that their price vector is
\[
P_c=\Omega'P,
\qquad
P=
\begin{pmatrix}
16\\
8\\
6
\end{pmatrix}.
\]
Thus,
\begin{align*}
P_c
&=
\begin{pmatrix}
0 & \frac{1}{10} & -\frac{1}{10} \\
-\frac{1}{10} & \frac{1}{10} & \frac{1}{5} \\
\frac{1}{10} & -\frac{1}{10} & -\frac{1}{10}
\end{pmatrix}
\begin{pmatrix}
16\\
8\\
6
\end{pmatrix} \\
&=
\begin{pmatrix}
\frac{1}{5}\\
\frac{2}{5}\\
\frac{1}{5}
\end{pmatrix}.
\end{align*}

\subsection*{(c)}

Asset $D$ has payoff vector
\[
X_D=
\begin{pmatrix}
30\\
20\\
10
\end{pmatrix}.
\]
By the law of one price,
\[
P_D=P_c'X_D
=\frac{1}{5}(30)+\frac{2}{5}(20)+\frac{1}{5}(10)
=16.
\]

\subsection*{(d)}

An SDF must price each contingent claim. For state $n$,
\[
P_c(n)=\pi(n)M(n),
\]
so
\[
M(n)=\frac{P_c(n)}{\pi(n)}.
\]
Therefore,
\[
M=
\begin{pmatrix}
\dfrac{1/5}{1/4}\\[4pt]
\dfrac{2/5}{1/2}\\[4pt]
\dfrac{1/5}{1/4}
\end{pmatrix}
=
\begin{pmatrix}
\frac{4}{5}\\
\frac{4}{5}\\
\frac{4}{5}
\end{pmatrix}.
\]
The SDF is strictly positive in every state. Hence, for any nonzero payoff
$X_A\geq 0$,
\[
P_A=\mathbb{E}[MX_A]>0,
\]
which rules out arbitrage. The SDF is also constant across states, so
$\operatorname{Cov}(M,X)=0$ for every payoff $X$. Thus, there is no
state-dependent risk premium coming from the SDF: a dollar of payoff is valued
the same across states after accounting for state probabilities.

\subsection*{(e)}

Using the SDF from part (d),
\begin{align*}
P_D
&=\mathbb{E}[MX_D] \\
&=\frac{4}{5}\left(\frac{1}{4}(30)+\frac{1}{2}(20)
    +\frac{1}{4}(10)\right) \\
&=16,
\end{align*}
which is the same price obtained in part (c).

\section*{Question 3(a)}

The investor solves
\[
\max_{w_t}\ \mathbb{E}_t[U(W_{t+1})]
\qquad\text{subject to}\qquad
W_{t+1}=W_t w_t'R_{t+1},
\quad
\mathbf{1}'w_t=1.
\]
The Lagrangian is
\[
\mathcal{L}
=\mathbb{E}_t\!\left[U\!\left(W_t w_t'R_{t+1}\right)\right]
+\lambda_t(1-\mathbf{1}'w_t).
\]
The first-order condition is
\[
\mathbb{E}_t\!\left[U'(W_{t+1})W_tR_{t+1}\right]
-\lambda_t\mathbf{1}=0.
\]
For any two assets $j$ and $k$, their first-order conditions imply
\[
\mathbb{E}_t\!\left[U'(W_{t+1})W_tR_{j,t+1}\right]
=
\mathbb{E}_t\!\left[U'(W_{t+1})W_tR_{k,t+1}\right].
\]
Subtracting gives
\[
\mathbb{E}_t\!\left[
U'(W_{t+1})W_t(R_{j,t+1}-R_{k,t+1})
\right]=0.
\]
Let asset $k$ be the baseline asset, so that
$xR_{j,t+1}=R_{j,t+1}-R_{f,t+1}$. Under CRRA utility,
\[
U'(W_{t+1})=W_{t+1}^{-\gamma}
=W_t^{-\gamma}R_{w,t+1}^{-\gamma}.
\]
Consequently,
\[
W_t^{1-\gamma}
\mathbb{E}_t\!\left[R_{w,t+1}^{-\gamma}xR_{j,t+1}\right]=0.
\]
Because $W_t^{1-\gamma}>0$, this becomes
\[
\mathbb{E}_t\!\left[R_{w,t+1}^{-\gamma}xR_{j,t+1}\right]=0
\qquad\text{for every }j.
\]
Therefore, the SDF
\[
M_{t+1}=R_{w,t+1}^{-\gamma}
\]
prices the vector of excess returns:
\[
\mathbb{E}_t\!\left[R_{w,t+1}^{-\gamma}xR_{t+1}\right]=0.
\]

\section*{Question 4(a)}

Let
\[
\mu_t=\mathbb{E}_t[R_{w,t+1}].
\]
Using a first-order Taylor expansion of
$M_{t+1}=R_{w,t+1}^{-\gamma}$ around $\mu_t$ gives
\begin{align*}
M_{t+1}
&\approx
\mu_t^{-\gamma}
-\gamma\mu_t^{-\gamma-1}(R_{w,t+1}-\mu_t) \\
&=(1+\gamma)\mu_t^{-\gamma}
-\gamma\mu_t^{-\gamma-1}R_{w,t+1}.
\end{align*}
Hence,
\[
M_{t+1}\approx a_t-b_tR_{w,t+1},
\]
where
\[
a_t=(1+\gamma)\mu_t^{-\gamma},
\qquad
b_t=\gamma\mu_t^{-\gamma-1}.
\]
One disadvantage of the linear SDF is that it can become negative for a
sufficiently large realization of $R_{w,t+1}$. By contrast, the original
power-utility SDF $R_{w,t+1}^{-\gamma}$ is always positive whenever
$R_{w,t+1}>0$.
